\section{VLM choice for embodied EQA on a single consumer GPU}
\label{app:model_choice}

Our full Habitat EQA stack --- simulator, voxel map, SigLIP grounding encoder, and the
multimodal LLM used for captioning, answering, and answer debiasing --- must share one
24\,GB consumer GPU. This appendix records the resulting model bake-off: which open
vision--language models we evaluated, how they were quantized to fit, and the perhaps
counterintuitive result that a \emph{grounding-specialized} VLM at 8B parameters
outperforms both the prior 3B control \emph{and} a newer 9B generalist that leads
Qwen3-VL on static VQA benchmarks~\citep{qwen35}.

\subsection{Roles and memory budget}

The VLM serves three roles in Dynagraph: (i) per-view captioning and scene-graph object
labeling, (ii) the iterative EQA answer/explore decision, and (iii) the answer-debiasing
votes (free-form answer matching with choice-rotation fallback). Grounding and
open-vocabulary retrieval are \emph{not} VLM roles: they require a contrastive dual
encoder, so SigLIP~\citep{zhai2023siglip} is retained regardless of VLM choice (loading
it in bf16 halves its footprint to $\sim$1.75\,GB with pixel-feature cosine $>$0.999 vs.\
fp32). The non-VLM stack (Habitat, voxel map, SigLIP) holds $\sim$7\,GB.
On an RTX~4090 during live Qwen3-VL-8B int4 HM-EQA we measure $\sim$13\,GiB used / $\sim$11\,GiB free, so 8--9B-class models sit comfortably under 4-bit NF4~\citep{dettmers2023qlora}.
A \textbf{32B int4} checkpoint (e.g.\ \texttt{Qwen3-VL-32B-Instruct}) is a \emph{candidate} with OOM risk rather than ruled out: peak VRAM must be measured with \texttt{scripts/smoke\_vl\_model.py} before any Habitat holdout (see \texttt{docs/habitat/vlm\_bakeoff.md}, ``Wave: larger VLM'').

\subsection{Candidates}

\begin{table}[h]
  \centering
  \caption{VLM candidates for the Dynagraph EQA bake-off (single RTX 24\,GB GPU).}
  \label{tab:vlm_candidates}
  \small
  \begin{tabular}{@{}lcccl@{}}
    \toprule
    Model & Params & Precision & Stack VRAM & Notes \\
    \midrule
    Qwen2.5-VL-3B~\citep{qwen25vl} & 3B & bf16 & $\sim$13\,GB & control; int4 unstable for this arch \\
    Qwen3-VL-8B~\citep{qwen3vl} & 8B & int4 & $\sim$13\,GB & native bbox grounding \\
    Qwen3-VL-8B~\citep{qwen3vl} & 8B & fp16 (Orin) & $\sim$18\,GB & same model, unquantized; served remotely \\
    Qwen3.5-9B~\citep{qwen35} & 9B & int4 & $\sim$13\,GB & natively multimodal; hybrid DeltaNet \\
    Gemma-4-E4B~\citep{gemma4} & 4B eff. & int4 & $\sim$12\,GB & MatFormer; no native grounding \\
    \bottomrule
  \end{tabular}
\end{table}

Two practical notes. First, Qwen3.5's hybrid linear-attention stack (3:1 Gated
DeltaNet:attention) silently falls back to slow PyTorch kernels unless the
\texttt{fla-core} Triton kernels are installed; without them single episodes exceeded a
4-hour budget, and with them per-episode time returned to the $\sim$30--50 minute range
typical of the 8B. Second, bitsandbytes int4 was numerically unstable for the
Qwen2.5-VL-3B architecture in our stack, so the control runs in bf16.

\subsection{Protocol and exclusions}

All candidates ran \textbf{Dynagraph} with the same exploration budget
(20 planning steps, frontier-node exploration, keyword weight 2) and the same
answer-debiasing pipeline (free-form match, then choice-rotation vote). We compared models
on a \textbf{canonical-6} HM-EQA subset (question ids 3, 14, 17, 28, 35, 81) after excluding
two ids that broke the harness rather than the models: id~31 never completed within
budget on int4 8--9B models (0/16 correct historically on the 3B control) and id~94
triggered a scene-graph blowup ($\sim$670 nodes vs.\ $\sim$70 typical) that starved the
EQA loop to zero iterations. The winner was promoted to a \textbf{balanced-31} subset (31 ids;
same two exclusions). Reproduction commands and result paths are documented in the project
repository under \texttt{docs/habitat/vlm\_bakeoff.md}.

\subsection{Results}

\begin{table}[h]
  \centering
  \caption{Canonical-6 HM-EQA accuracy (Dynagraph + debiasing). Gold answer letters:
  q3=B, q14=D, q17=D, q28=D, q35=D, q81=D.}
  \label{tab:vlm_bakeoff}
  \small
  \begin{tabular}{@{}lcc@{}}
    \toprule
    Model & Correct & $\Delta$ vs.\ 3B control \\
    \midrule
    \textbf{Qwen3-VL-8B int4} & \textbf{5/6 (83\%)} & \textbf{+3} \\
    Qwen3.5-9B int4 & 3/6 (50\%) & +1 \\
    Qwen2.5-VL-3B bf16 (control) & 2/6 (33\%) & --- \\
    Gemma-4-E4B int4 & 2/5 & +0 (q14 OOM; phase skipped) \\
    \bottomrule
  \end{tabular}
\end{table}

\begin{table}[h]
  \centering
  \caption{Per-question outcomes on canonical-6 (+ = correct, -- = wrong, n/a = episode
  not completed). Both Qwen models miss q17 (exploration never observed the target object).}
  \label{tab:vlm_per_question}
  \small
  \begin{tabular}{@{}lcccccc@{}}
    \toprule
    Model & q3 & q14 & q17 & q28 & q35 & q81 \\
    \midrule
    Qwen3-VL-8B & + & + & -- & + & + & + \\
    Qwen3.5-9B & + & + & -- & -- & + & -- \\
    Qwen2.5-VL-3B & -- & + & -- & -- & -- & + \\
    Gemma-4-E4B & + & n/a & -- & -- & -- & + \\
    \bottomrule
  \end{tabular}
\end{table}

Upsizing from 3B to 8B on the \emph{same} pipeline and questions recovered three
additional correct answers (+50\% relative), establishing that the prior 3B ceiling was
model-limited rather than exploration-limited on this slice. The balanced-31 run with the
8B winner was in progress at submission time; we report its JSONL path in the reproduction
doc above.

\subsection{Quantization is not free: int4 vs.\ fp16 on the same 8B}
\label{app:model_choice_quantization}

The bake-off above is run on a single 24\,GB GPU, which forces 4-bit NF4
quantization~\citep{dettmers2023qlora} for the 8B winner. To isolate the cost of that
precision, we served the \emph{same} \texttt{Qwen3-VL-8B} unquantized in fp16 on a
separate NVIDIA AGX Orin (61\,GiB unified memory) over an OpenAI-compatible HTTP
endpoint (\texttt{EMET\_VL\_ENDPOINT}), and re-ran a representative 30-question subset of
HM-EQA with the identical Dynagraph harness, exploration budget, and debiasing pipeline.
Only the weight precision differs.

\begin{table}[h]
  \centering
  \caption{30-qid HM-EQA subset, same harness, same qids: fp16 vs.\ int4 (Qwen3-VL-8B).}
  \label{tab:vlm_precision}
  \small
  \begin{tabular}{@{}lcc@{}}
    \toprule
    Precision & Correct & Accuracy \\
    \midrule
    int4 (local 4090) & 9/30 & 30\% \\
    \textbf{fp16 (remote Orin)} & \textbf{16/30} & \textbf{53\%} \\
    \midrule
    $\Delta$ & \textbf{+7} & \textbf{+23\,pp} \\
    \bottomrule
  \end{tabular}
\end{table}

fp16 recovers seven additional correct answers (+78\% relative). The gain is not uniform:
it is concentrated in the count/clock slice (10/15 vs.\ 6/15), where the task demands
reading fine visual detail (clock hands, shelf counts, mat placement) that 4-bit
quantization distorts; on a separate 15-qid diverse slice fp16 was roughly flat (6/15
vs.\ 5/15). This is the first clean same-model measurement that quantization cost is
\emph{task-dependent} and material: an unquantized (or larger) VLM is a first-class lever,
not a footnote. Prior work reporting higher embodied-EQA accuracy with stronger VLMs may
partly reflect this precision effect rather than model-family superiority alone.

Operationally, the fp16 arm ran $\sim$3--4$\times$ slower (Orin nvgpu + HTTP JPEG
transport, $\sim$10--13\,min/qid vs.\ $\sim$3\,min int4 in-process), so precision is
traded against wall-clock on a fixed GPU; the remote-Orin split amortizes that by
offloading the VLM off the single 24\,GB card. Pitfalls for reproducing the remote arm
(fp16-only serve, \texttt{HF\_HOME}, ssh detachment, float32 OOM) are documented in
\texttt{docs/experiments/vlm\_quantization.md}.

\subsection{Why a newer, larger generalist loses}
\label{app:model_choice_failures}

Qwen3.5-9B is trained from scratch on interleaved image/video/text tokens and reports
substantially stronger \emph{static} VQA and reasoning scores than Qwen3-VL at similar
scale~\citep{qwen35}. On embodied MCQ EQA it still trailed the 8B specialist and barely
beat the 3B control. The gap is not explained by parameter count alone (9B $>$ 8B) nor by
aggregate benchmark rank; it reflects task-specific behaviors that HM-EQA amplifies:

\paragraph{Benchmarks vs.\ embodied EQA.}
Static VQA rewards fluent scene description and multi-step reasoning over a fixed image
set. Our task adds (a) \emph{coverage} --- the agent must navigate to see the evidence;
(b) \emph{literal counting and presence} --- ``how many red pillows'' and ``is the child
sleeping'' admit no hedging; and (c) \emph{format compliance} --- our debiasing stage
asks for a terse free-form answer without captions or option letters. Qwen3.5 excels at
the first skill (long ``Caption: \ldots'' preambles on every debias call) but loses on
the latter two.

\paragraph{Failure modes observed on canonical-6.}
\begin{enumerate}
  \item \textbf{Confident visual hallucination (q81).} Qwen3.5 asserted that ``image~1
  clearly shows a child sleeping on a red car-shaped bed'' and graded itself confident;
  no person appears in any view. Qwen3-VL answered ``No'' on the same scene. This is the
  largest single-question swing between the two Qwen models.
  \item \textbf{Perception / counting (q28).} Qwen3.5's free-form reply described ``one
  red pillow'' (gold: two); Qwen3-VL's reply stated ``two red pillows'' and matched the
  correct choice. Both had similar exploration budgets; the difference is literal visual
  aggregation, not MCQ letter bias.
  \item \textbf{Debiasing path disabled by instruction drift.} Qwen3.5 prefaced debias
  queries with ``Caption: \ldots'' despite explicit ``do not caption'' instructions, so
  the free-form matcher rarely engaged and answers fell through to rotation voting ---
  where Qwen3.5 showed \emph{position-locked} selection (votes A,B,C,D across four
  rotations of one question), the same failure mode we previously measured on the 3B
  control.
  \item \textbf{Shared coverage failure (q17).} Neither Qwen model observed the woven
  basket; both answered incorrectly. No answering strategy recovers missing evidence.
\end{enumerate}

Gemma-4-E4B underperformed similarly (2/5 completed episodes) and could not finish q14:
deterministic CUDA OOM immediately after load (VRAM fully exhausted on a 12\,MiB
allocation), whereas both Qwen models completed q14 without incident --- suggesting
family-specific memory spikes under our Habitat+SigLIP co-residency, not a universal
quantization limit.

\paragraph{Infrastructure interaction.}
The bake-off also surfaced a Dynagraph failure mode unrelated to model choice: certain
scenes inflate the scene graph by an order of magnitude (669--885 nodes vs.\ $\sim$70
typical), starving the main EQA loop. On q35 the \emph{debias vote alone} returned the
correct answer after the main loop produced zero iterations --- an unplanned robustness
benefit of the two-path answering design, and a reminder that prompt size from graph
memory can dominate VLM choice independently of model quality.

\subsection{Takeaways}

\begin{enumerate}
  \item \textbf{Prefer grounding-trained VLMs} (Qwen3-VL) over higher-benchmark
  generalists (Qwen3.5) at equal VRAM when the task requires literal presence/counting
  and tight instruction following.
  \item \textbf{$\sim$8B int4 is the practical ceiling} on a single 24\,GB GPU shared with
  Habitat and SigLIP; upsizing from 3B bf16 to 8B int4 yielded the largest accuracy gain
  we measured (+3/6 on canonical-6). But precision is the next lever: the same 8B in fp16
  (served on a remote Orin) adds another +7/30 (+23\,pp) on a 30-qid subset, concentrated
  on fine-detail count/clock questions.
  \item \textbf{Measure position bias and format compliance directly} on your MCQ pipeline
  --- they moved end-task accuracy more than headline benchmark deltas between Qwen
  generations.
  \item \textbf{Install \texttt{fla-core}} for Qwen3.5 if you keep it in the loop; otherwise
  wall-clock dominates comparisons unfairly.
\end{enumerate}
